% =============================================================================
% Хавсралт Г: Ажлын түүх
% =============================================================================

\chapter{Ажлын түүх}
\label{app:history}

Төсөл git-д бүртгэгдээгүй тул ажлын түүхийг \tsql{docs/WORKLOG.md}-ээс нэгтгэв (2026 он, Улаанбаатарын цаг).

\begin{longtable}{p{2.0cm}p{12.0cm}}
\caption{Ажлын гол үе шат}\label{tab:history}\\
\toprule
Огноо & Хийсэн ажил \\
\midrule
\endfirsthead
\toprule
Огноо & Хийсэн ажил \\
\midrule
\endhead
10-15 & Заавар ба туршилтын төлөвлөгөөг код болгосон: CLI, дамжлага, RAG, үнэлгээ, сургалт, аюулгүй байдал; dev 300, test 200, 160 жишээ \\
10-16 & Agentic горим, LaTeX тайлангийн анхны хувилбар; Kaggle ба Modal-д бэлтгэсэн \\
10-17 & A, B бүлэг (Kaggle T4); C1--C3 сургалт (Modal); D2 сургалт ба H бүлэг (Vast A100) \\
10-18 & E2 dev 41.0\%; test D2 21.0\%, H4 34.0\%, E2 38.5\%, E3c 46.5\%; 195 шинэ жишээ; E4 \\
10-19 & I3, I3t dev ба test; E4 test 45.0\%; D3 QLoRA 112/116 алхам; 280 жишээ \\
10-20 & T2 орчуулга (маск 122 → 37); agentic K1--K3 тохиргоо \\
10-23 & Орчуулгын fine-tune ба тогтвортой сургалтын төлөвлөгөө; MT v0 өгөгдөл \\
10-25 & MT v2 (2 828 хос), S1--S4, T-2/T-4; Kaggle STAGE=mt эхэлсэн; Text2SQL Studio; локал PostgreSQL демо \\
10-26 & Танилцуулга ба энэ тайлан; Kaggle дээр MADLAD-400-3B-ийн QLoRA сургалт (Version 13) \\
\bottomrule
\end{longtable}
